\section{Events as evidence rather than prose}\label{sec:C13:events}
An event record states that a particular observation or transition was recorded under specified conditions. It should distinguish a requested action from a completed effect, and a model's prediction from an environment observation. A proposed minimal envelope includes an event identifier, causal parents, producer, observed time, effective time where known, payload type, schema version, and integrity metadata. Sensitive payloads can be stored separately behind access controls while the envelope retains a reference and retention status.

The distinction between observed and effective time is essential for corrections. A user can report today that a preference changed last week. A system can discover a failed external operation after it previously recorded a timeout. Ordering by ingestion time alone does not establish which belief is currently valid. The event log should preserve both when information entered the system and what interval it describes. A query for the latest known fact and a query for what the agent knew at a past decision then become different operations.

Let $e_i\prec e_j$ denote a causal dependency. A tool result depends on its invocation, and a summary depends on the source events it abstracts. Causal order need not coincide with a single globally synchronized clock. Lamport's treatment of event ordering supplies the foundational distinction between causal precedence and a compatible total ordering~\cite{lamport1978clocks}. For agent logs, the practical consequence is to retain causal links rather than assuming that wall-clock sorting reconstructs the decision process.

Typed events enable different validation rules. A model proposal can be accepted as a proposal without being accepted as a fact. A tool receipt can require a validated tool identity and matching invocation identifier. A completion event can require specific tests. An authorization event can require an authorized principal and a scope. If all four are serialized as untyped text, later summarization can accidentally promote the first into the last. Data typing reduces this risk but cannot determine whether the content of a tool observation is itself truthful.

An append-only evidence log does not imply indefinite retention of all raw payloads. Deletion, redaction, and access revocation can be represented by policy-controlled transformations while preserving the fact that an earlier decision referenced an object that is no longer accessible. The system should avoid claiming that a deleted payload remains reconstructible if its retention policy intentionally removes it. Auditability and privacy create a real design tradeoff; an implementation must state which forms of evidence survive.

The resulting event history is different from a polished narrative. It can contain conflicting reports, partial failures, and superseded plans. Those are necessary evidence for diagnosing why the agent acted as it did. A model-facing summary is a derived view of that history, with its own version and provenance, rather than a replacement for the record of events.

